\textbf{Most of the reported performance of just-in-time software defect
prediction is an artefact of how it is measured.} On the corpus studied
here, a result that reads as 0.41 MCC under the field's standard
protocol is 0.10 under a defensible one --- and the single largest
component of that gap is circular scoring: evaluating models against the
same heuristic labels that trained them.

That is possible because JIT-SDP labels almost universally descend from
the SZZ algorithm, which infers defect-introducing commits by blaming
the lines a bug fix modifies. The inference fails whenever bug-fixing
commits are tangled, and because no independent ground truth is
available, models are conventionally scored against the very labels
whose errors they may have learned.

\textbf{This report covers the first two phases of that measurement},
across 21 Apache Java projects and 27,319 commits, holding out an
independently constructed reference label set so that a model trained on
SZZ can be scored against something other than SZZ.

\textbf{Phase 1 --- label quality.} Six SZZ variants are scored over an
identical denominator. No variant exceeds \textbf{27.2\% precision}.
Noise is class-conditional rather than symmetric (false-alarm rates
6.7--26.3\%, miss rates 35.9--73.3\%). Variants agree with one another
at \ensuremath{\kappa} up to \textbf{0.933} while agreeing with the reference at \ensuremath{\kappa}
\textbf{0.158--0.202} --- they share failure modes, so mutual agreement
is not evidence of correctness. Decomposing the misses shows that for
the most aggressive variants, \textbf{more than half of what their
filters remove is correct}: refinement has relocated error rather than
reduced it.

\textbf{Phase 2 --- downstream impact.} Three models across seven label
sources and four evaluation regimes, 16,380 runs, tested by paired
Wilcoxon at the project level with Hodges--Lehmann estimates, bootstrap
intervals and Holm correction applied both within family and globally.
Random \emph{k}-fold cross-validation inflates a high-capacity model by
\textbf{+0.137 MCC} (21 of 21 projects); scoring against the training
labels inflates it by a further \textbf{+0.230}. The two protocol
effects together are roughly \textbf{nine times} the label-source
effect, and both scale with model capacity. Under honest streaming
evaluation with reconstructed verification latency, training on the
reference labels beats all six SZZ variants, five of six surviving
global correction.

A result reading \textbf{0.41 MCC} under the field's standard
configuration is \textbf{0.10} under a defensible one. Because
self-scoring is only possible where labels are heuristic, the labelling
failure and the evaluation failure are not independent: the first
creates the second.

\textbf{Scope.} Two questions are raised here and deliberately left
open: which of SZZ's two error types an online learner cannot survive,
and whether a learner can be built to withstand it. Chapter 6 states
both designs, together with the specific controls this report has
already shown they will require, and records in advance the commitment
that the mitigation phase will be pre-registered.
